\clearpage
\VerdierPhysicalPageAnchor{0112}%
\begingroup
\small
\setlength{\abovedisplayskip}{4pt plus 2pt minus 2pt}%
\setlength{\belowdisplayskip}{4pt plus 2pt minus 2pt}%

\VerdierSemanticAnchor{ch2-1-1-3}
\noindent\textbf{定義 1.1.3.}\quad
$\mathcal D$ と $\mathcal D'$ を二つの三角圏とする。$\mathcal D$ から
$\mathcal D'$ への\emph{完全関手}とは、（必ずしも厳密ではない）
$\mathbb Z$-関手（第 I 章、1.2.4）
$(F,p)\colon\mathcal D\longrightarrow\mathcal D'$ であって、
$\mathcal D$ の任意の完全三角

\VerdierSemanticAnchor{ch2-1-1-3-source-triangle}
\[
  X\xrightarrow{u}Y\xrightarrow{v}Z\xrightarrow{w}X[1]
\]
に対し、三角

\VerdierSemanticAnchor{ch2-1-1-3-image-triangle}
\[
  FX\xrightarrow{Fu}FY\xrightarrow{Fv}FZ
  \xrightarrow{p(X)Fw}(FX)[1]
\]
が完全三角となるものをいう (1.1.4.2)。

\VerdierSemanticAnchor{ch2-1-1-4}
\medskip
\noindent\textbf{定義 1.1.4.}\quad
$\mathcal D_i$（$1\leq i\leq n$）および $\mathcal D$ を合わせて
$n+1$ 個の三角圏とする。$\prod_i\mathcal D_i$ から $\mathcal D$ への
\emph{多重完全関手}とは、捩れた $\mathbb Z^n$-関手（第 I 章、1.4.4）
$(F,p)\colon\prod_i\mathcal D_i\longrightarrow\mathcal D$ であって、
それぞれ圏 $\mathcal D_i$ に属する任意の対象族 $X_i$（$i\ne j$）に対し、
$\mathbb Z$-関手

\VerdierSemanticAnchor{ch2-1-1-4-partial-functor}
\[
  X\longmapsto F(X_1,\ldots,X,\ldots,X_n)\colon
  \mathcal D_j\longrightarrow\mathcal D
\]
が完全関手となるものをいう。

\VerdierSemanticAnchor{ch2-1-1-4-1}
\medskip
\noindent\textbf{1.1.4.1.}\quad
定義 (1.1.4) では、圏 $\prod_i\mathcal D_i$ にその標準的な
$\mathbb Z^n$-構造を入れ、群 $\mathbb Z^n$ は準同型

\VerdierSemanticAnchor{ch2-1-1-4-1-group-action}
\[
  \mathbb Z^n\longrightarrow\mathbb Z,
\]
すなわち座標の和を通じて圏 $\mathcal D$ に作用する（第 I 章、1.6.6）。

\VerdierSemanticAnchor{ch2-1-1-4-2}
\medskip
\noindent\textbf{1.1.4.2.}\quad
符号規則 I（第 I 章、1.6.6）に従い、特に断らない限り、多重完全関手は
$\mathbb Z^n$ の標準基底 $e_1,\ldots,e_n$ に関する Koszul 余サイクルで
捩るものとする（第 I 章、1.6.4）。

\VerdierSemanticAnchor{ch2-1-1-4-2-trivial-cocycle}
とりわけ、これは $\mathcal D$ から $\mathcal D'$ への完全関手が自明な
余サイクルで捩られることを意味する。

\VerdierSemanticAnchor{ch2-1-1-5}
\medskip
\noindent\textbf{定義 1.1.5.}\quad
$\mathcal D$ を三角圏、$\mathcal A$ をアーベル圏とする。関手
$H\colon\mathcal D\longrightarrow\mathcal A$ が\emph{コホモロジー的関手}
であるとは、任意の完全三角

\VerdierSemanticAnchor{ch2-1-1-5-source-triangle}
\[
  X\xrightarrow{u}Y\xrightarrow{v}Z\xrightarrow{w}X[1],
\]
に対し、図式

\VerdierSemanticAnchor{ch2-1-1-5-image-sequence}
\[
  HX\xrightarrow{Hu}HY\xrightarrow{Hv}HZ\xrightarrow{Hw}H(X[1])
\]
が $\mathcal A$ の完全列となることをいう。
\endgroup
